% Part: normal-modal-logic
% Chapter: filtrations
% Section: S5-decidable

\documentclass[../../../include/open-logic-section]{subfiles}

\begin{document}

\olfileid{nml}{fil}{dec}

\olsection{\Log{S5} د پرېکړې وړ دے}
د متناهي مدل خاصيت موږ ته د دې ښودلو اسانه لار راکوي چې د سکيماوو
له مخې ورکړل شوي مودالي منطقي نظامونه \emph{د پرېکړې وړ} دي (يعنې
داسې محاسبه‌کېدونکې کړنلاره شته چې وټاکي آيا کوم !!a{formula} په
نظام کښې !!{derivable} دے که نه)۔

\begin{thm}
  \Log{S5} د پرېکړې وړ دے۔
\end{thm}

\begin{proof}
  $!A$ راکړل شوے وګڼئ، او فرض کړئ چې په~$!A$ کښې راغلي
  !!{propositional variable}s له $p_1$, \dots, $p_k$ څخه وي۔ ځکه
  د هر $n$ لپاره د~$n$ نړۍو او $p_1$, \dots, $p_k$ ته د
  ارزښت ټاکلو لرونکي يوازې متناهي ډېر مدلونه شته، موږ په
  \emph{موازي ډول} دوه لړۍ شمېرلے شو: له يوې خوا په يوې منظمه
  طريقه د ثبوتونو په جوړولو د~\Log{S5} ټولې قضيې، او له بلې خوا
  هغه ټول عمومي مدلونه چې $1$, $2$, \dots نړۍې لري، او بيا
  وګورو چې آيا $!A$ په کوم داسې مدل کښې په يوه نړۍ ناسم دے۔
  په منجمده سرچينه کښې د دې دويمې پلټنې لپاره د عمومي مدل قيد
  نۀ دے ويل شوے؛ د همدې نظام د ضدمدل لپاره همدا قيد اړين دے۔
  له دغو دوو موازي بهيرونو څخه يو به بالاخره ځواب ورکړي، ځکه د
  \olref[com][fra]{thm:generaldet} او
  \olref[fmp]{cor:S5fmp} له مخې يا $!A$ !!{derivable} دے، يا
  په يوه متناهي عمومي مدل کښې ناسم دے۔
\end{proof}

پورتنے ثبوت د~\Log{S5} لپاره ځکه کار کوي چې د عمومي مدلونو
فلټريشنونه په خپله عمومي وي۔ د انعکاسيت او مسلسل‌والي لپاره هم
همداسې ده، خو د نورو خاصيتونو لپاره زيات کار پکار دے۔

\end{document}
